\section{Introduction}

The capacity of a female mammal to enter puberty and sustain adult reproduction depends on whether body energy stores are adequate to underwrite the metabolic cost of gestation and lactation. The adipocyte-derived hormone leptin communicates this information to the brain and acts as a permissive metabolic signal for the reproductive axis: mice and humans lacking leptin or the leptin receptor (LepR) are obese and fail to enter puberty, and exogenous leptin reverses hypogonadotropic hypogonadism in leptin-deficient subjects~\cite{ahima1997leptin,farooqi2014twenty}. Leptin acts through its receptor (LepRb) on neurons distributed across many hypothalamic and brainstem nuclei, yet the cellular identity and neurochemical effectors through which the signal reaches gonadotropin-releasing hormone (GnRH) neurons have remained incompletely resolved~\cite{louis2011molecular}.

GnRH neurons do not themselves express LepRb, implying that leptin acts on intermediary populations that then gate GnRH output. Candidate relays include Kiss1/neurokinin B/dynorphin (KNDy) neurons of the arcuate nucleus (Arc) and proopiomelanocortin (POMC) neurons of the Arc, as well as LepRb neurons of the ventral premammillary nucleus (PMv). Genetic deletion of LepRb specifically from Kiss1 neurons does not disrupt puberty or fertility, and only a small minority ($\sim$10\%) of Arc KNDy neurons coexpress LepRb, arguing against a direct kisspeptin-centred model of leptin action~\cite{donato2011leptin,cravo2011kiss1,louis2011molecular,true2011leptin}. A melanocortin relay in which Arc POMC neurons signal onto KNDy neurons has been proposed to convey leptin's permissive effect on puberty, consistent with direct LepRb action on POMC cells~\cite{manfredilozano2016defining,berglund2012direct}. In contrast, earlier work demonstrated that bilateral lesions of the PMv blunt leptin-induced sexual maturation and that unilateral endogenous re-expression of LepRb in the PMv alone restores puberty and fertility in LepR-null females, pinpointing the PMv as a privileged relay~\cite{donato2011leptin,mahany2018obesity}.

PMv LepRb neurons show the signatures expected of a metabolic-to-reproductive relay: leptin depolarizes about three quarters of them through a PI3K/TRPC-dependent mechanism, they monosynaptically innervate GnRH neurons, and they are activated by both leptin and sexual odorants~\cite{williams2011acute,leshan2009direct}. Importantly, the majority of PMv LepRb neurons are glutamatergic (Vglut2$^+$), raising the largely untested hypothesis that glutamate is the neurotransmitter through which leptin's permissive signal is conveyed to the reproductive axis. Contribution of glutamate to this circuit has been inferred from transmitter-specific LepR deletions elsewhere in the brain: selective removal of LepRb from GABAergic, but not glutamatergic, populations produces severe obesity and reproductive defects, suggesting that the cell-type requirements for body-weight regulation and reproductive regulation differ and may segregate across transmitter phenotypes~\cite{martin2014leptinresponsive}. In parallel, Kiss1 neurons of the Arc are themselves predominantly glutamatergic, and glutamatergic signalling has been implicated in the reciprocal KNDy circuit that shapes pulsatile GnRH release~\cite{cravo2011kiss1,navarro2013interactions,navarro2011regulation}.

Despite this convergent evidence, no study has tested causally whether activation of PMv LepRb neurons is sufficient to drive GnRH/luteinizing hormone (LH) release, whether conditional deletion of Vglut2 from LepRb cells impairs puberty and adult reproduction, or whether the requirement for glutamate is specifically local to the PMv. Addressing these questions is essential both for understanding how body energy stores gate reproduction and for rationalising the apparently conflicting evidence on the role of kisspeptin and melanocortin circuits as leptin relays~\cite{louis2011molecular,amstalden2011neuroendocrine,gottsch2004role,iqbal2001evidence}.

Here we address these gaps by combining three complementary loss- and gain-of-function strategies in mice. First, we use chemogenetic stimulation to test whether acute activation of PMv LepRb neurons in fed adult females is sufficient to trigger LH release on a physiologically relevant timescale. Second, we delete Vglut2 from all LepRb cells to ask whether glutamate released from LepRb neurons is required for pubertal maturation, estrous cyclicity, pituitary responsiveness to kisspeptin-10 and ultimately fertility. Third, we combine selective LepR re-expression in the PMv with concomitant deletion of Vglut2 in the same population of LepR-null mice, to ask whether glutamate is specifically the permissive neurotransmitter at the PMv node. Together, these experiments define glutamate as an indispensable effector through which PMv LepRb neurons relay the metabolic signal of leptin to the hypothalamic--pituitary--gonadal axis.

\section{Related Work}

\paragraph{Leptin, energy balance and reproductive permission.} Leptin links adiposity to fertility in mice and humans: ob/ob and db/db mice and humans with congenital leptin or leptin-receptor mutations are obese and hypogonadal, and leptin replacement restores reproductive competence~\cite{ahima1997leptin,farooqi2014twenty}. This framing of leptin as a permissive rather than instructive signal motivates the search for leptin-responsive neurons whose discrete activation is sufficient to drive gonadotropin release, independently of body weight change.

\paragraph{Candidate relays between LepRb and GnRH neurons.} Because GnRH neurons themselves do not express LepRb, leptin's action must be transmitted through intermediary neurons~\cite{louis2011molecular}. Genetic removal of LepRb from Kiss1 neurons does not disrupt puberty or fertility~\cite{donato2011leptin}, and only a small minority of Arc KNDy neurons coexpress LepRb~\cite{cravo2011kiss1,true2011leptin}. A melanocortin pathway, in which POMC neurons integrate leptin signals and act via MC3/4R on Kiss1 neurons, has been proposed based on pharmacology and virogenetics~\cite{manfredilozano2016defining}. Deletion of LepRb from GABAergic but not glutamatergic populations produces obesity and delayed puberty, implying that the transmitter phenotype of LepRb neurons matters for reproductive function~\cite{martin2014leptinresponsive}. Our work asks the complementary question at the circuit level, focusing on glutamate at the PMv.

\paragraph{The PMv as a reproductive-metabolic node.} The PMv was identified as a key leptin-responsive site for reproduction by lesion experiments and by endogenous re-expression studies in LepR-null mice~\cite{donato2011leptin,mahany2018obesity}. PMv LepRb neurons are monosynaptically connected to GnRH neurons and are activated by leptin and by sexual odorants~\cite{leshan2009direct}. Electrophysiological recordings show that leptin depolarises a majority of PMv LepRb neurons through PI3K/TRPC signalling~\cite{williams2011acute}. However, the transmitter identity through which PMv LepRb neurons exert their reproductive action has not been directly tested.

\paragraph{Kisspeptin, NKB and the KNDy network.} Kisspeptin, signalling through Kiss1r, is a direct and indispensable stimulator of GnRH neurons; exogenous kisspeptin-10 elicits robust LH release in mice~\cite{gottsch2004role}. Within the Arc, kisspeptin is co-expressed with neurokinin B (Tac2) and dynorphin in KNDy neurons, forming an autofeedback circuit that shapes pulsatile GnRH release~\cite{navarro2011regulation,navarro2013interactions}. This network is downstream of metabolic inputs, and Tac2/NKB expression is sensitive to energetic status~\cite{amstalden2011neuroendocrine}. Our kisspeptin-10 challenge and mediobasal hypothalamic qPCR are designed to probe where in this chain glutamatergic PMv inputs act.

\paragraph{Cell-type-specific dissection of hypothalamic circuits.} Conditional Cre/lox deletion of Lepr or Vglut2, AAV-Cre re-expression in Lepr-reactivatable alleles, and DREADD-based chemogenetic activation together enable circuit dissection with cell-type and regional resolution~\cite{berglund2012direct,kim2012foxo1}. This study combines all three strategies to separate the necessity and sufficiency of PMv glutamatergic output from the broader effects of leptin on body weight, food intake and generic LepRb signalling.
